\section{从教学 ISA 过渡到真实指令编码}

教学 ISA 用少量规则说明“操作码控制数据通路”，真实 ISA 则必须让数十年的软件、编译器和处理器实现共同遵守稳定契约。这个契约不仅包含加减和跳转，还包括寄存器宽度、地址计算、异常优先级、原子操作、内存序、特权级与扩展发现机制。学习真实 ISA 时，应把架构规定的软件可见结果与某颗处理器内部的流水线实现分开。

以 RISC-V 的基础整数指令为例，常见指令长度为 32 bit，操作码位于固定低位，寄存器编号也位于较固定的位置，使译码和寄存器读取容易并行。立即数为了服务 I、S、B、U、J 等格式被分散编码，硬件译码后再拼接并符号扩展。分支立即数最低位不编码，是因为目标地址按至少 2 B 对齐；省下的位用于扩大可达范围。

x86 则保留变长编码和大量历史语义。一条指令可能由前缀、操作码、ModR/M、SIB、位移和立即数组成。前端先识别边界、查表译码，再把复杂指令拆成内部微操作。两类 ISA 的表面编码差异很大，但进入执行后仍要解决寄存器依赖、地址生成、访存权限、异常与退休等共同问题。

\begin{longtable}{L{0.18\linewidth}L{0.31\linewidth}L{0.41\linewidth}}
\toprule
观察层次 & 需要回答的问题 & 不应混淆的对象 \\
\midrule
指令编码 & 哪些 bit 表示操作、寄存器和立即数 & 汇编器展示的助记符可能是伪指令 \\
ISA 语义 & 执行后哪些架构状态发生变化 & 内部可能拆成多个微操作 \\
ABI & 参数、返回值、保存寄存器如何约定 & ABI 是软件约定，不是每条指令的硬件语义 \\
微架构 & 何时发射、旁路、预测和退休 & 不得改变 ISA 承诺的可见结果 \\
\bottomrule
\end{longtable}

\section{地址生成、访存宽度与端序}

load/store 指令通常用基址寄存器加符号扩展位移形成有效地址。随后地址要经过对齐检查、页表翻译、权限检查和 Cache 访问。访问宽度决定读取或写入几个字节；较窄的 load 还必须明确零扩展或符号扩展。端序规定多字节数值与逐字节地址的对应关系，不改变寄存器内部 bit 的编号。

未对齐访问能否完成由 ISA 和实现共同约束。有的 ISA 允许硬件拆成两次访问，有的环境报告对齐异常，由软件模拟。若一次访问跨越页面或设备寄存器边界，不能假定“两边都成功或都失败”；体系结构必须规定异常的精确性，驱动也应避免对 MMIO 做未声明的宽度与对齐组合。

原子读改写不能简单理解为 load 后接 store，因为其间可能被另一核心插入访问。实现可使用独占监视、Cache line 所有权或互连锁定来保证原子提交。compare-and-swap 还把读取值与比较结果返回软件，使无锁算法能够检测竞争并重试。

\section{特权级、CSR 与同步陷阱}

普通程序不能直接修改页表基址、关闭中断或访问任意设备寄存器。处理器用特权级和控制状态寄存器划分权限；系统调用指令主动进入内核，非法指令、缺页和除零等同步异常则由当前指令触发。外部中断与执行中的指令异步，但处理器仍选择一个可恢复的架构边界进入处理程序。

陷阱入口至少要保存原因、故障地址或补充状态、返回地址与先前特权状态，再跳到向量入口。处理程序修复原因后，可从原指令重试，也可跳过指令或终止任务。重试型异常要求故障指令尚未提交部分可见副作用；这正是乱序核心必须通过按序退休维护精确异常的原因。

\section{内存序与屏障不是 Cache 刷新}

编译器和处理器都可能重排互不依赖的内存操作。单核程序若看不出差异，这类优化通常合法；设备寄存器和多核共享数据却会观察到先后次序。内存模型规定普通 load/store、原子操作和屏障之间允许哪些结果。release 操作约束它之前的写先于发布动作可见，acquire 操作约束观察发布结果后的访问不能越过它提前完成。

屏障约束顺序和完成边界，不必然把所有 Cache 内容写到 DRAM。Cache clean/invalidate 处理副本与所有权，持久化 flush 处理掉电边界，三者目的不同。向非一致性 DMA 设备提交描述符时，软件通常先写内存，再做所需的 Cache 维护与写屏障，最后写 doorbell；完成方向则先确认设备完成并处理 Cache 可见性，再读取结果。

\section{扩展、能力发现与兼容性}

现代 ISA 通过扩展增加向量、加密、虚拟化或位操作能力。软件不能只根据处理器品牌猜测功能，而应读取体系结构规定的能力位或由操作系统提供的能力接口。二进制若直接执行不存在的扩展，会触发非法指令异常；库通常用启动时分派，在通用实现和加速实现之间选择。

兼容性还包括状态保存。新增向量寄存器意味着上下文切换、信号处理和调试器必须知道其格式；为降低代价，操作系统可按“首次使用”懒惰启用，但必须防止一个任务读取另一个任务留下的状态。可迁移虚拟机通常只暴露各目标主机都支持的能力子集，否则迁移后无法继续执行。

\section*{章末练习}

\begin{longtable}{L{0.08\linewidth}L{0.32\linewidth}L{0.5\linewidth}}
\toprule
题号 & 类型 & 题目 \\
\midrule
1 & 层次辨析 & 分别说明 ISA、ABI 和微架构规定什么，并举一个只属于各层的例子。 \\
2 & 地址计算 & 某 load 用寄存器基址加负位移访问 16-bit 数据。列出从立即数译码到寄存器得到结果之间的检查和变换。 \\
3 & 异常 & 缺页异常为何通常重试原指令，而系统调用返回时为何通常从下一条指令继续？ \\
4 & 原子性 & 解释普通 load 加 store 为何不能替代 compare-and-swap。 \\
5 & 内存序 & 描述 CPU 向非一致性 DMA 设备发布描述符时，数据写、Cache 维护、屏障与 doorbell 的合理顺序。 \\
6 & 兼容性 & 程序怎样安全使用一个并非所有处理器都支持的向量扩展？还需哪些操作系统支持？ \\
\bottomrule
\end{longtable}

\section*{参考解答与要点}

\begin{longtable}{L{0.08\linewidth}L{0.32\linewidth}L{0.5\linewidth}}
\toprule
题号 & 类型 & 解答要点 \\
\midrule
1 & 层次辨析 & ISA 规定指令和架构状态，如加法结果；ABI 规定调用接口，如参数寄存器；微架构规定内部实现，如是否拆成微操作和流水级数。 \\
2 & 地址计算 & 拼接并符号扩展立即数，与基址相加；检查对齐；进行地址翻译和权限检查；读取两个字节；按端序组合；按指令语义零扩展或符号扩展。 \\
3 & 异常 & 缺页处理修复映射后需再次执行尚未完成的访存；系统调用指令本身已完成“请求进入内核”的作用，返回地址通常指向其后继。 \\
4 & 原子性 & 两条普通指令之间可被其他核心访问，比较条件可能失效；原子指令把读、条件检查和写作为不可分割的提交事务。 \\
5 & 内存序 & 先填描述符，执行所需 clean/同步，再用写屏障保证数据先可见，最后写 doorbell。完成后还需按平台规则失效或同步 Cache 再读结果。 \\
6 & 兼容性 & 通过标准能力接口检测扩展并运行时分派，保留通用后备路径；操作系统还须启用、保存和恢复新增寄存器状态。 \\
\bottomrule
\end{longtable}
